% ----------------------------------------------------------------------------------------------------------------------
% AI Almanach - Released under the MIT License
% ----------------------------------------------------------------------------------------------------------------------


% User guides
% Latex: https://ftp.agdsn.de/pub/mirrors/latex/dante/macros/latex/base/usrguide.pdf
% Tabularray: https://mirror.physik.tu-berlin.de/pub/CTAN/macros/latex/contrib/tabularray/tabularray.pdf
% TColorbox: https://mirror.clientvps.com/CTAN/macros/latex/contrib/tcolorbox/tcolorbox.pdf
% Posterbox Tutorial: https://mirror.clientvps.com/CTAN/macros/latex/contrib/tcolorbox/tcolorbox-tutorial-poster.pdf#[0,{%22name%22:%22Fit%22}]
% Xcolor: https://mirror.clientvps.com/CTAN/macros/latex/contrib/xcolor/xcolor.pdf
% TikZ and PGFPlots: https://tikz.dev/ and https://tikz.dev/pgfplots/
%                    https://tug.ctan.org/info/visualtikz/VisualTikZ.pdf 
% SVG Package: https://ctan.org/pkg/svg?lang=en
% Geometry: https://ftp.gwdg.de/pub/ctan/macros/latex/contrib/geometry/geometry.pdf
% fontspec: https://texdoc.org/serve/fontspec/0
% fancyhdr: https://ftp.fau.de/ctan/macros/latex/contrib/fancyhdr/fancyhdr.pdf
% biblatex: https://distrib-coffee.ipsl.jussieu.fr/pub/mirrors/ctan/macros/latex/contrib/biblatex/doc/biblatex.pdf
% amsmath: https://www.latex-project.org/help/documentation/amsldoc.pdf
% grapheur: https://ctan.tetaneutral.net/graphics/pgf/contrib/tkz-grapheur/doc/tkz-grapheur-doc-en.pdf
% enumitem: https://mirror.ibcp.fr/pub/CTAN/macros/latex/contrib/enumitem/enumitem.pdf

\documentclass[8pt]{extreport}
\tracinglostchars=2

% The managed build starts from the repository root.
% Search document-local files without depending on the caller's directory.
\makeatletter
\def\input@path{{ai almanach/}{./}}
\makeatother

% --- Hyphenation Rules --------------------------------------------------------
\usepackage[USenglish]{babel}

% --- Fonts --------------------------------------------------------------------

% Load system fonts via fontspec (requires XeLaTeX or LuaLaTeX)
% to get a list of installed fonts: luaotfload-tool --list=format 
\usepackage{fontspec}
\IfFontExistsTF{FreeSans}{
    \setmainfont{FreeSans}[Scale=1]
    \setsansfont{FreeSans}[Scale=1]
    \setmonofont{FreeMono}[Scale=1]
}{
    \setmainfont{Latin Modern Roman}
    \setsansfont{Latin Modern Sans}
    \setmonofont{Latin Modern Mono}
}

% Prefer the poster font when installed, but keep the book reproducible.
\IfFontExistsTF{QTFuturePoster}{
    \newfontfamily\headerfont{QTFuturePoster}
}{
    \newfontfamily\headerfont{Latin Modern Sans}
}


% Emoji support
\usepackage{emoji}

% --- Page layout --------------------------------------------------------------
\usepackage[
    % showframe, 
    nomarginpar, 
    a4paper,
    portrait, 
    right=1.1cm,
    left=1.1cm,
    top=1.35cm,
    bottom=1.15cm,
    headsep=0.3cm,
    headheight=0.75cm,
    footskip=0.6cm,
]{geometry}


\usepackage{fancyhdr}
\pagestyle{fancy}


% --- TColorbox ---------------------------------------------------------------
\usepackage{tcolorbox}
\tcbuselibrary{skins,raster,listings,breakable,minted, poster}


% --- Tables -------------------------------------------------------------------
\usepackage{tabularray}

% --- Color Names -------------------------------------------------------------
\usepackage[dvipsnames]{xcolor}

% --- Graphics -----------------------------------------------------------------
\usepackage{graphicx}
\usepackage{svg}
\graphicspath{{./figures/}{./icons/}{./logos/}}
\usepackage{tikz}
\usepackage{pgfplots}
\pgfplotsset{compat=1.18}
\usetikzlibrary{shapes.geometric}    % for shapes
\usetikzlibrary{arrows.meta}         % for arrows, Arrow tip library 
\usetikzlibrary{positioning}         % for relative positioning
\usetikzlibrary{calc}                % for calculations to make complex coordinate calculations
%\usetikzlibrary{backgrounds}        % Background Library "defines background for pictures". To use this in a Tikzpicture, an option is passed, e.g. \begin{tikzpicture}[show background rectangle], with a background rectangle style defined before the picture. (e.g. \tikzset{background rectangle/.style={<define background rectangle style here>}}
% \usetikzlibrary{calendars}         % for calendar drawings
% \usetikzlibrary{fadings}           % for fading effects
% \usetikzlibrary{patterns}          % for patterns
% \usetikzlibrary{shadows}           % for shadow effects 
\usetikzlibrary{chains}            % for chain effects
\usetikzlibrary{fit}               % for fitting nodes together
% \usetikzlibrary{er}                % for entity relationship diagrams
% \usetikzlibrary{intersections}     % for calculating intersections of paths
% \usetikzlibrary{mindmap}           % for mind maps
\usetikzlibrary{matrix}            % for matrix drawings
\usetikzlibrary{angles}           % for angle markings
\usetikzlibrary{quotes}           % for quotes in angles
% \usetikzlibrary{trees}             % for tree drawings
% \usetikzlibrary{decorations.pathmorphing}  % for path decorations like zigzag lines
\usetikzlibrary{decorations.pathreplacing} % for braces and other path decorations
% \usetikzlibrary{decorations.markings}      % for markings along a path
% \usetikzlibrary{decorations.shapes}        % for shape decorations
% \usetikzlibrary{decorations.text}          % for text decorations along a path
% \usetikzlibrary{pgfplots.colorbrewer}      % for color maps

% \usepackage{forest}               % for drawing trees https://ctan.ceremade.dauphine.fr/graphics/pgf/contrib/forest/forest-doc.pdf

% \usepackage{adjustbox}          % for adjusting boxes

\usepackage[labelfont=sf]{caption}  % for figure captions on non-floats

\input{shapes/flowchart_styles.tex}  % load flowchart styles


% --- Math Environments --------------------------------------------------------
\usepackage{amsmath}
\usepackage{amssymb}
\usepackage{nicematrix}
\NiceMatrixOptions{cell-space-limits = 1pt}
\usepackage{siunitx}

% --- Bibliography -------------------------------------------------------------
\usepackage[autostyle=true]{csquotes} % recommended before biblatex
\usepackage[backend=biber,style=numeric,sorting=nyt]{biblatex}
\addbibresource{bibliography/references.bib}

\renewcommand{\bibfont}{\normalfont\footnotesize}


% --- Enumerations -----------------------------------------------------------------
\usepackage{enumitem}

\setlist{noitemsep}
\setlist[1]{labelindent=\parindent} % < Usually a good idea
\setlist[itemize]{leftmargin=*}
\setlist[itemize,1]{label=-}

% --- Helpers -----------------------------------------------------------------
\usepackage{multicol}
\usepackage{microtype}

% --- URL, cross-references, and PDF metadata ----------------------------------
\colorlet{citecolor}{black}
\colorlet{linkcolor}{black}
\colorlet{urlcolor}{black}
\usepackage[
  bookmarks=true,
  breaklinks=true,
  pdfborder={0 0 0},
  citecolor=citecolor,
  linkcolor=linkcolor,
  urlcolor=urlcolor,
  colorlinks=true,
  linktocpage=false,
  hyperindex=true,
  linkbordercolor=white]{hyperref}
\usepackage[capitalise,nameinlink]{cleveref}
\hypersetup{
    pdftitle={AI Almanach: Machine Learning, Data Science, and Intelligent Systems for Engineers},
    pdfauthor={Bernhard Gerlach},
    pdfsubject={A didactic engineering reference for artificial intelligence and data science},
    pdfkeywords={artificial intelligence, machine learning, data science, deep learning, large language models, evaluation}
}


% --- Options ---------------------------------------------------------
% \pagestyle{empty}
% \setlength\parindent{0pt}
% \setlength{\tabcolsep}{2pt}
% \baselineskip=0pt
\setlength{\columnsep}{0.65cm}
% \setlength{\parskip}{0.1\baselineskip}

\setlist[enumerate]{itemsep=0.2em, topsep=0.2em, parsep=0em, partopsep=0em}

\captionsetup{font=small,labelfont={bf,sf}}
% \captionsetup[sub]{font=small,labelfont={bf,sf}}

\setcounter{secnumdepth}{4}


\newcommand{\maincolor}{LimeGreen!90!Gray}
\newcommand{\maincolspan}{4}
\newcommand{\graphscale}{0.9}

\title{AI Almanach}
\author{Bernhard Gerlach}
\date{\today}

% --- TColorbox Skins ----------------------------------------------------------

% Colorbox skin for bash listings
\tcbsubskin{mintedbash}{empty}{
        size=minimal,
        listing engine=minted,
        listing only,
        minted style=friendly,
        minted language=bash,
        minted options={
            fontsize=\footnotesize,
            breaklines,
            autogobble,
        },
        colback=gray!10!white,
        colframe=gray!10!white,
        listing only,
        left=0em,
        enhanced,
}

% Colorbox skin for python listings
\tcbsubskin{mintedpython}{empty}{
        size=minimal,
        listing engine=minted,
        listing only,
        minted style=friendly,
        minted language=python,
        minted options={
            fontsize=\footnotesize,
            breaklines,
            autogobble,
        },
        colback=gray!10!white,
        colframe=gray!10!white,
        listing only,
        left=0em,
        enhanced,
        before skip=0.3\baselineskip,
}

% Skin for header box
\tcbsubskin{headerboxskin}{empty}{
        size=minimal,
        coltitle=black!10!black,
}



% Skin for orange-based section boxes
\tcbsubskin{sectionboxskin}{standard}{
    size=small,
    left=0.55em,
    right=0.55em,
    top=0.45em,
    bottom=0.45em,
    boxrule=0.45pt,
    arc=0.8mm,
    colback=\maincolor!2!white,
    colframe=\maincolor!75!black,
    coltitle=\maincolor!10!black,
    colbacktitle=\maincolor!75!white,
    fonttitle=\sffamily\bfseries,
    toptitle=0.35mm,
    bottomtitle=0.35mm,
    fontupper=\sffamily\small,
    fontlower=\sffamily\small,
    lower separated=false,
    valign lower=center,
    halign=left,
    subtitle style={top=0.4mm, bottom=0mm, boxrule=0.2pt, boxsep=0.1mm,
        colback=\maincolor!50!\maincolor!20!white,
        colupper=\maincolor!50!gray,
        fontupper=\sffamily\bfseries\small,
        coltext=\maincolor!10!black,
        height=1.1\baselineskip,
      }, 
}


% Skin for sub-section boxes
\tcbsubskin{subsectionboxskin}{standard}{
    size=fbox, 
    top=0.4em,
    % left=0.2em,
    % right=0.2em,
    % arc=0.25mm,
    % top=0.4mm, 
    % bottom=0mm, 
    % boxrule=0.2pt, 
    % boxsep=0.1mm,
    frame empty,
    % colback=\maincolor!2!white,
    colback=white,
    colbacktitle=\maincolor!50!\maincolor!20!white,
    fonttitle=\sffamily\bfseries\small,
    coltitle=\maincolor!10!black,
    % height=1.1\baselineskip,    
}


% Skin for simple tables inside boxes
\tcbsubskin{boxtablesimple}{empty}{
    size=minimal,
    before skip=0.2\baselineskip,
    after skip=0.2\baselineskip,
    colback=white,
    colframe=gray,
    frame empty,
}


\tcbsubskin{sectionraster}{sectionboxskin}{
    raster columns=6,
    raster row skip=0.35cm,
    raster column skip=0.35cm,
    skin=sectionboxskin,
}

% --- Semantic learning boxes -------------------------------------------------
% These environments standardise how new material is introduced throughout v2.
\newtcolorbox{learningobjectives}[1][Learning objectives]{
    enhanced,
    breakable,
    title={#1},
    colback=RoyalBlue!3!white,
    colframe=RoyalBlue!65!black,
    fonttitle=\sffamily\bfseries,
    fontupper=\sffamily\small,
    boxrule=0.45pt,
    arc=0.6mm
}

\newtcolorbox{plainlanguage}[1][In plain language]{
    enhanced,
    breakable,
    title={#1},
    colback=LimeGreen!4!white,
    colframe=LimeGreen!55!black,
    fonttitle=\sffamily\bfseries,
    fontupper=\sffamily\small,
    boxrule=0.45pt,
    arc=0.6mm
}

\newtcolorbox{engineernote}[1][Engineer's note]{
    enhanced,
    breakable,
    title={#1},
    colback=Orange!5!white,
    colframe=Orange!70!black,
    fonttitle=\sffamily\bfseries,
    fontupper=\sffamily\small,
    boxrule=0.45pt,
    arc=0.6mm
}

\newtcolorbox{workedexample}[1][Worked example]{
    enhanced,
    breakable,
    title={#1},
    colback=Cyan!3!white,
    colframe=Cyan!55!black,
    fonttitle=\sffamily\bfseries,
    fontupper=\sffamily\small,
    boxrule=0.45pt,
    arc=0.6mm
}

\newtcolorbox{checkpoint}[1][Check your understanding]{
    enhanced,
    breakable,
    title={#1},
    colback=Violet!3!white,
    colframe=Violet!55!black,
    fonttitle=\sffamily\bfseries,
    fontupper=\sffamily\small,
    boxrule=0.45pt,
    arc=0.6mm
}

\newtcolorbox{practicallab}[1][Practical lab]{
    enhanced,
    breakable,
    title={#1},
    colback=TealBlue!3!white,
    colframe=TealBlue!65!black,
    fonttitle=\sffamily\bfseries,
    fontupper=\sffamily\small,
    boxrule=0.6pt,
    arc=0.6mm
}

\newtcolorbox{cautionbox}[1][Common trap]{
    enhanced,
    breakable,
    title={#1},
    colback=Red!3!white,
    colframe=Red!60!black,
    fonttitle=\sffamily\bfseries,
    fontupper=\sffamily\small,
    boxrule=0.45pt,
    arc=0.6mm
}

\newcommand{\keyterm}[1]{\textbf{\textcolor{\maincolor!65!black}{#1}}}
\newcommand{\evidencelevel}[1]{\textsf{\footnotesize Evidence: #1}}

% --- Chapter 7-inspired learning rhythm --------------------------------------
% Every chapter begins with an advance organiser and ends with retrieval and
% transfer prompts. The visual sits beside the explanatory text so a short
% introduction does not leave an empty half-page.
\newcommand{\chapteropening}[6]{%
\begin{tcolorbox}[
    enhanced,
    colback=white,
    colframe=\maincolor!70!black,
    boxrule=0.65pt,
    arc=1mm,
    title={Chapter compass},
    fonttitle=\sffamily\bfseries,
    before skip=0.5\baselineskip,
    after skip=0.7\baselineskip]
\begin{minipage}[t]{0.49\linewidth}
\sffamily
\textbf{Essential question}\par
#2

\medskip
\textbf{Learning outcome}\par
Explain the chain from \emph{#3} to \emph{#6}, calculate or trace one small
case, diagnose a failure at each boundary, and name evidence that would make
an engineering decision defensible.

\medskip
\textbf{How to read}\par
First build the mental model, then follow the formal mechanism, and finally
challenge the assumptions with the worked examples and retrieval prompts.
\end{minipage}\hfill
\begin{minipage}[t]{0.46\linewidth}
\centering
\begin{tikzpicture}[
    node distance=3.5mm,
    roadmap step/.style={
        draw=\maincolor!70!black,
        fill=\maincolor!7!white,
        rounded corners=0.8mm,
        minimum width=0.82\linewidth,
        minimum height=6.5mm,
        align=center,
        font=\sffamily\scriptsize},
    roadmap arrow/.style={-{Stealth[length=1.8mm]},thick,
                          draw=\maincolor!70!black}]
\node[roadmap step] (a) {#3};
\node[roadmap step,below=of a] (b) {#4};
\node[roadmap step,below=of b] (c) {#5};
\node[roadmap step,below=of c] (d) {#6};
\draw[roadmap arrow] (a) -- (b);
\draw[roadmap arrow] (b) -- (c);
\draw[roadmap arrow] (c) -- (d);
\end{tikzpicture}
\captionof{figure}{Chapter roadmap: #3, #4, #5, and #6.}
\label{fig:#1-roadmap}
\end{minipage}
\end{tcolorbox}
}

\newcommand{\chaptersynthesis}[6]{%
\section{Synthesis and Retrieval Practice}
\begin{tcbitemize}[skin=sectionraster]
\tcbitem[title={Connect the mechanism},raster multicolumn=3]
Start with \textbf{#3} and trace what information, uncertainty, or control
signal crosses each boundary until \textbf{#6}.
For every transition, state the representation, the operation, and the
assumption that makes the operation meaningful.

\tcbitem[title={Failure-first review},raster multicolumn=3]
Choose one plausible failure at each stage.
Separate specification error from data error, model error, implementation
error, and operating-context change.
Then identify the earliest observable signal that could reveal each failure.

\tcbitem[title={Evidence ladder},raster multicolumn=6]
Use the checklist as a design-review record rather than a vocabulary test.
\tcblower
\begin{tblr}{
    width=\linewidth,
    colspec={Q[l,0.18\linewidth]X[l]X[l]},
    row{1}={font=\bfseries,bg=\maincolor!15!white},
    hlines,
    vlines}
Stage & Deep-dive question & Minimum useful evidence \\
#3 & What is represented, measured, or assumed? & Provenance, units, bounds,
and a counterexample \\
#4 & Which mechanism transforms the input, and why? & A hand-worked case,
invariant, or executable test \\
#5 & How is uncertainty or error exposed? & Baseline comparison, slices,
stress tests, and uncertainty \\
#6 & Which decision follows, and who owns the residual risk? & Acceptance
criterion, monitoring signal, fallback, and owner \\
\end{tblr}
\captionof{table}{Deep-dive evidence checklist for this chapter.}
\label{tab:#1-synthesis-checklist}

\tcbitem[title={Retrieve, explain, transfer},raster multicolumn=6]
Without looking back, answer the chapter question: \emph{#2}
Sketch the #3--#6 chain from memory, explain one equation or algorithm in
plain language, and adapt the chain to a new domain.
Check your answer against the chapter, correct it in a different colour, and
write one question that the chapter does not settle.
\end{tcbitemize}
}

\makeatletter
\def\almanachchapterkey{}
\def\almanachchapterquestion{}
\def\almanachchapterstagea{}
\def\almanachchapterstageb{}
\def\almanachchapterstagec{}
\def\almanachchapterstaged{}
\let\almanachoriginalchapter\chapter
\RenewDocumentCommand{\chapter}{s o m}{%
    \IfBooleanTF{#1}{%
        \almanachoriginalchapter*{#3}%
    }{%
        \IfNoValueTF{#2}{%
            \almanachoriginalchapter{#3}%
        }{%
            \almanachoriginalchapter[#2]{#3}%
        }%
        \ifx\almanachchapterkey\@empty
        \else
            \chapteropening
                {\almanachchapterkey}
                {\almanachchapterquestion}
                {\almanachchapterstagea}
                {\almanachchapterstageb}
                {\almanachchapterstagec}
                {\almanachchapterstaged}%
        \fi
    }%
}
\makeatother

\newcommand{\almanachinputchapter}[7]{%
    \gdef\almanachchapterkey{#1}%
    \gdef\almanachchapterquestion{#2}%
    \gdef\almanachchapterstagea{#3}%
    \gdef\almanachchapterstageb{#4}%
    \gdef\almanachchapterstagec{#5}%
    \gdef\almanachchapterstaged{#6}%
    \input{#7}%
    \almanachdeepdive{#1}%
    \chaptersynthesis{#1}{#2}{#3}{#4}{#5}{#6}%
    \gdef\almanachchapterkey{}%
}

% Chapter-specific extensions injected by \almanachinputchapter.
% Keep each extension compact, current, and tied to an engineering decision.

\expandafter\def\csname almanachdeepdive@orientation\endcsname{%
\section{Deep Dive: The Foundation-Model Lifecycle}
\begin{tcbitemize}[skin=sectionraster]
\tcbitem[title={One Model, Many System Layers},raster multicolumn=3]
A foundation model is only one component in a lifecycle that includes data
rights, pretraining, post-training, evaluation, retrieval, tools, serving,
monitoring, and incident response.
The same weights can behave very differently under another prompt template,
retrieval corpus, decoding policy, tool permission, or user workflow.

\tcbitem[title={Capability Is Not a Deployment Decision},raster multicolumn=3]
Translate a proposed capability into a user, decision, error cost, operating
boundary, baseline, acceptance test, fallback, and owner.
Multimodal input or fluent generation expands the interface; it does not remove
the need to define what counts as correct and safe.

\tcbitem[title={Mini Design Review},raster multicolumn=6]
For one AI application, draw
\emph{data $\rightarrow$ model $\rightarrow$ system $\rightarrow$ human
$\rightarrow$ outcome}.
At every arrow, name what can be lost, distorted, attacked, or misunderstood.
Then choose one metric and one qualitative investigation that could reveal the
failure before deployment.
\end{tcbitemize}
}

\expandafter\def\csname almanachdeepdive@mathematics\endcsname{%
\section{Deep Dive: Numerical Stability and Honest Uncertainty}
\begin{tcbitemize}[skin=sectionraster]
\tcbitem[title={Conditioning Before Optimisation},raster multicolumn=3]
An algorithm can be mathematically correct and numerically unreliable.
Inspect scale, units, rank, and condition number before interpreting a large
gradient or unstable coefficient.
Standardisation, orthogonal decompositions, and regularisation often improve
the problem being solved rather than merely speeding up a solver.

\tcbitem[title={Log-Sum-Exp},raster multicolumn=3]
Directly computing $\log\sum_i e^{x_i}$ may overflow.
With $m=\max_i x_i$,
\begin{equation}
  \log\sum_i e^{x_i}=m+\log\sum_i e^{x_i-m}.
\end{equation}
Subtracting $m$ changes neither the result nor softmax ratios, while keeping
exponents in a representable range.

\tcbitem[title={Intervals Need Assumptions},raster multicolumn=6]
A confidence, credible, bootstrap, prediction, or conformal interval answers a
different question under different assumptions.
For every interval, state the random object, conditioning information,
coverage target, data-dependence assumptions, and whether tuning reused the
calibration data.
Check empirical coverage on relevant slices rather than trusting the nominal
percentage by name.
\end{tcbitemize}
}

\expandafter\def\csname almanachdeepdive@programming\endcsname{%
\section{Deep Dive: Reproducible ML Software}
\begin{tcbitemize}[skin=sectionraster]
\tcbitem[title={Notebook to Package},raster multicolumn=3]
Keep exploration in notebooks, but move reusable transformations, models, and
metrics into typed modules with documented interfaces.
A command-line entry point plus a versioned configuration makes the experiment
rerunnable without executing cells in a remembered order.

\tcbitem[title={Test Properties, Not Only Examples},raster multicolumn=3]
Example tests check known cases.
Property and metamorphic tests check invariants: shuffling rows should not
change a permutation-invariant aggregate; serialization should preserve a
prediction within tolerance; a monotone transform should preserve ranks.
These tests catch families of failures that one golden input misses.

\tcbitem[title={Operational Skeleton},raster multicolumn=6]
Build a minimal repository with locked dependencies, configuration validation,
structured logging, data-schema checks, unit tests, one end-to-end smoke test,
an environment capture, and a single command that trains and evaluates a
baseline.
Record code revision, data identifier, seed, hardware, and output artifact in
the same run manifest.
\end{tcbitemize}
}

\expandafter\def\csname almanachdeepdive@data-analytics\endcsname{%
\section{Deep Dive: Data Contracts, Lineage, and Drift}
\begin{tcbitemize}[skin=sectionraster]
\tcbitem[title={A Schema Is Not a Data Contract},raster multicolumn=3]
A contract includes semantics, units, allowed values, null meaning, freshness,
uniqueness, privacy classification, owners, change policy, and service level.
Two columns can share a type while representing incompatible populations or
measurement procedures.

\tcbitem[title={Lineage Makes Errors Reversible},raster multicolumn=3]
Record source, transformation version, validation result, and downstream
consumers for every material dataset and feature.
Lineage should answer which models used a corrupted source and which reports
must be recomputed after a correction.

\tcbitem[title={Drift Is a Question, Not One Score},raster multicolumn=6]
Distinguish input drift $p(x)$, label drift $p(y)$, concept drift $p(y\mid x)$,
and operational changes in collection or policy.
A PSI or KS alert can indicate distribution change but cannot identify cause or
business impact by itself.
Pair statistical signals with slice-level performance, data-quality checks,
and an investigation playbook.
\end{tcbitemize}
}

\expandafter\def\csname almanachdeepdive@machine-learning\endcsname{%
\section{Deep Dive: From Score to Decision Policy}
\begin{tcbitemize}[skin=sectionraster]
\tcbitem[title={Probability, Threshold, Action},raster multicolumn=3]
A classifier score becomes useful only through a decision policy.
Choose thresholds from false-positive and false-negative consequences,
capacity constraints, calibration, and an abstention or review path.
Accuracy at the default threshold rarely captures this operating decision.

\tcbitem[title={Calibration and Coverage},raster multicolumn=3]
Measure reliability curves and proper scoring rules on held-out data.
When uncertainty is material, compare calibrated probabilities, prediction
intervals, or conformal sets and report empirical coverage and set size by
slice.
Coverage without useful sharpness can still produce an unusable system.

\tcbitem[title={Data-Centric Error Loop},raster multicolumn=6]
After a baseline, cluster errors by mechanism: ambiguous target, label error,
missing feature, distribution gap, shortcut, capacity limit, or threshold
policy.
Make one intervention at a time, preserve a protected evaluation set, and
measure whether improvement transfers beyond the error slice that inspired the
change.
\end{tcbitemize}
}

\expandafter\def\csname almanachdeepdive@deep-learning\endcsname{%
\section{Deep Dive: Modern Training and Inference}
\begin{tcbitemize}[skin=sectionraster]
\tcbitem[title={Residual Pre-Norm Blocks},raster multicolumn=3]
Modern deep networks stabilise optimisation with residual paths and
normalisation.
In a pre-norm Transformer block,
$x' = x + \mathrm{Attention}(\mathrm{Norm}(x))$ and
$y = x' + \mathrm{MLP}(\mathrm{Norm}(x'))$.
The residual path carries information and gradients even when a learned branch
is initially poor.

\tcbitem[title={Training Parallelism vs Inference},raster multicolumn=3]
Self-attention can process training positions in parallel, but dense attention
still has quadratic interaction cost in sequence length.
Autoregressive generation remains sequential across new tokens.
A key--value cache avoids recomputing earlier keys and values, trading memory
for lower per-token latency.

\tcbitem[title={Efficiency Is Multi-Dimensional},raster multicolumn=6]
Quantisation, distillation, sparsity, low-rank adaptation, activation
checkpointing, batching, and distributed sharding reduce different costs and
introduce different errors.
Benchmark end-to-end quality, latency, throughput, peak memory, energy, and
operational complexity under the intended workload; parameter count alone is
not a deployment metric.
\end{tcbitemize}
}

\expandafter\def\csname almanachdeepdive@evaluation-mlops\endcsname{%
\section{Deep Dive: Evaluating Generative and Adaptive Systems}
\begin{tcbitemize}[skin=sectionraster]
\tcbitem[title={Quality Is a Vector},raster multicolumn=3]
For generative systems, separate task success, factual grounding, instruction
following, safety, style, latency, and cost.
Aggregate scores can hide a release-blocking failure, so define non-compensable
gates for critical behaviours.

\tcbitem[title={Judges Need Evaluation Too},raster multicolumn=3]
An LLM judge is a measurement instrument.
Calibrate it against blinded human decisions, randomise order, test position
and verbosity bias, report agreement and uncertainty, and keep a stable anchor
set across judge changes.

\tcbitem[title={Sequential Release Evidence},raster multicolumn=6]
Combine offline suites with shadow traffic, canaries, guarded online
experiments, rollback triggers, and post-release monitoring.
Account for repeated looks at online metrics and delayed outcomes.
A release is an evidence update, not permission to retire the evaluation set.
\end{tcbitemize}
}

\expandafter\def\csname almanachdeepdive@reinforcement-learning\endcsname{%
\section{Deep Dive: Offline, Safe, and Preference-Based Learning}
\begin{tcbitemize}[skin=sectionraster]
\tcbitem[title={Offline RL and Coverage},raster multicolumn=3]
Offline RL learns from a fixed behaviour dataset.
Actions outside that dataset's support can receive optimistic values because
their outcomes are weakly identified.
Report behaviour-policy coverage and compare with behaviour cloning and a
conservative non-learning baseline.

\tcbitem[title={Safety Is a Constraint},raster multicolumn=3]
Reward penalties do not guarantee a hard safety limit.
Constrained policies, shields, action filters, simulators, and independent
supervisors address different layers.
Evaluate rare violations, recovery, and reward hacking rather than only mean
return.

\tcbitem[title={Off-Policy Evidence},raster multicolumn=6]
Before deploying a new policy, estimate its value with importance weighting,
doubly robust estimators, or a validated simulator.
Large policy shifts create high-variance or biased estimates when action
probabilities have little overlap.
Report effective sample size, sensitivity to clipping, and where evaluation
depends on extrapolation.
\end{tcbitemize}
}

\expandafter\def\csname almanachdeepdive@trustworthy-ai\endcsname{%
\section{Deep Dive: Agentic Systems and Deterministic Boundaries}
\begin{tcbitemize}[skin=sectionraster]
\tcbitem[title={Treat Tool Output as Untrusted},raster multicolumn=3]
Retrieved documents, web pages, messages, and tool results can contain
instructions designed to compete with system policy.
Parse data with schemas, isolate it from instructions, and keep authorization
outside the model.

\tcbitem[title={Least-Privilege Tool Use},raster multicolumn=3]
Issue short-lived credentials for the minimum resource and action.
Validate arguments, enforce limits, require confirmation for consequential
changes, and log the request, decision, tool result, and artifact versions.
The model proposes; a deterministic boundary disposes.

\tcbitem[title={Trajectory-Level Assurance},raster multicolumn=6]
Evaluate complete trajectories, including retries, tool selection, state
changes, recovery, and human handover.
A component that passes isolated prompts can still fail through compounding
errors, excessive autonomy, or a permission path that bypasses the intended
control.
\end{tcbitemize}
}

\expandafter\def\csname almanachdeepdive@computer-vision\endcsname{%
\section{Deep Dive: Foundation Vision and Robust Perception}
\begin{tcbitemize}[skin=sectionraster]
\tcbitem[title={Prompts Change the Vision Interface},raster multicolumn=3]
Promptable segmentation and vision--language models replace a fixed class list
with points, boxes, text, or examples.
This flexibility expands evaluation: test localisation, open-vocabulary
recognition, grounding, refusal, and consistency across prompt wording.

\tcbitem[title={Shift Is Often Physical},raster multicolumn=3]
Weather, optics, compression, sensor ageing, viewpoint, motion blur, and scene
composition alter the image-generating process.
Corruption benchmarks are useful probes but do not replace data collected in
the intended operational domain.

\tcbitem[title={Perception-to-Action Trace},raster multicolumn=6]
For one detected object, trace calibration, preprocessing, model output,
tracking, fusion, uncertainty, planning, and fallback.
Express how spatial or temporal error propagates into an action margin.
Evaluate the system-level consequence, not only IoU or mAP in isolation.
\end{tcbitemize}
}

\expandafter\def\csname almanachdeepdive@nlp-generative-ai\endcsname{%
\section{Deep Dive: Retrieval, Tools, and Structured Generation}
\begin{tcbitemize}[skin=sectionraster]
\tcbitem[title={RAG Has Two Evaluations},raster multicolumn=3]
Evaluate retrieval with relevance, recall, ranking, freshness, permissions, and
coverage.
Evaluate generation with grounding, citation correctness, completeness, and
task success.
A faithful answer to irrelevant evidence is still a retrieval failure.

\tcbitem[title={Structured Output Is a Contract},raster multicolumn=3]
Define a schema, validate every field, reject unknown actions, and separate
syntactic validity from semantic authorization.
Repair loops need bounded retries and telemetry; silently accepting a plausible
but invalid object moves uncertainty into downstream code.

\tcbitem[title={Long Context Is Not Perfect Memory},raster multicolumn=6]
More tokens increase cost and can dilute relevant evidence.
Compare long-context prompting with retrieval, summarisation, and explicit
state.
Test information at different positions, conflicting evidence, multilingual
inputs, and instructions embedded inside retrieved content.
\end{tcbitemize}
}

\expandafter\def\csname almanachdeepdive@practical-llm-engineering\endcsname{%
\section{Deep Dive: Serving the Adapted Model}
\begin{tcbitemize}[skin=sectionraster]
\tcbitem[title={Quantisation Experiment},raster multicolumn=3]
Compare the reference model with 8-bit and 4-bit weight variants on the same
quality, latency, throughput, and peak-memory suite.
Separate quantised loading from quantisation-aware training and record which
layers or tensors remain at higher precision.

\tcbitem[title={Measure the Workload Shape},raster multicolumn=3]
Prompt length, generated length, batch concurrency, cache reuse, and hardware
determine serving behaviour.
Report time to first token and inter-token latency in addition to total
throughput; an average can hide an unusable tail.

\tcbitem[title={Adapter Governance},raster multicolumn=6]
Version base model, tokenizer, chat template, adapter, merge procedure,
training data, licence, and evaluation result as one deployable bill of
materials.
Verify a merged artifact against the separate adapter path before promotion and
retain a rollback target.
\end{tcbitemize}
}

\expandafter\def\csname almanachdeepdive@fintech\endcsname{%
\section{Deep Dive: Point-in-Time Financial Evaluation}
\begin{tcbitemize}[skin=sectionraster]
\tcbitem[title={What Was Knowable Then?},raster multicolumn=3]
Reconstruct features, universe membership, prices, and filings as they were
available at the decision timestamp.
Later corrections, surviving firms, or revised labels create look-ahead and
survivorship bias even when rows are split by date.

\tcbitem[title={Backtest the Executable Policy},raster multicolumn=3]
Include transaction costs, spread, latency, market impact, turnover,
constraints, rejected orders, and capital limits.
Compare against a simple implementable strategy and test performance across
regimes rather than selecting a convenient interval.

\tcbitem[title={Decisions Need Reasons and Appeals},raster multicolumn=6]
For credit, fraud, or AML workflows, connect model output to threshold,
capacity, adverse-action explanation, human review, and appeal.
Measure calibration and error costs by relevant groups, plus queue burden and
time to resolution.
\end{tcbitemize}
}

\expandafter\def\csname almanachdeepdive@healthcare\endcsname{%
\section{Deep Dive: From Retrospective Accuracy to Clinical Utility}
\begin{tcbitemize}[skin=sectionraster]
\tcbitem[title={External and Prospective Evidence},raster multicolumn=3]
Validate across sites, devices, time periods, workflows, and patient groups.
A retrospective external set tests transportability; a prospective silent
study tests current data flow; an impact study tests whether use changes care
and outcomes.

\tcbitem[title={Thresholds Depend on the Pathway},raster multicolumn=3]
Sensitivity and specificity do not determine value without prevalence,
capacity, downstream tests, delay, and harm.
Decision-curve or utility analysis makes the intervention threshold explicit
and should be complemented by calibration and subgroup uncertainty.

\tcbitem[title={Change Control for Learning Systems},raster multicolumn=6]
Define which data, preprocessing, model, threshold, interface, or workflow
changes require revalidation.
Monitor input and outcome delay, override patterns, incidents, and calibration.
A safe update plan includes rollback and communication, not only a newer model
file.
\end{tcbitemize}
}

\expandafter\def\csname almanachdeepdive@industrial-ai\endcsname{%
\section{Deep Dive: Digital Twins, Edge AI, and Safe Updates}
\begin{tcbitemize}[skin=sectionraster]
\tcbitem[title={Calibrate the Twin},raster multicolumn=3]
A digital twin is useful only within a validated envelope.
Estimate parameters from measured behaviour, compare residuals across operating
regimes, quantify uncertainty, and identify effects absent from the simulator.

\tcbitem[title={Edge Constraints Shape the Model},raster multicolumn=3]
Timing jitter, memory, power, connectivity, sensor bandwidth, and thermal limits
can dominate model choice.
Benchmark on the target device with realistic preprocessing and concurrent
loads, not only on a workstation.

\tcbitem[title={Safe Over-the-Air Change},raster multicolumn=6]
Sign artifacts, bind model and configuration versions, stage rollout, verify a
shadow or canary, protect compatibility, and keep an independently tested
fallback.
An update must preserve the operational safety envelope even when its average
prediction metric improves.
\end{tcbitemize}
}

\expandafter\def\csname almanachdeepdive@supply-chain\endcsname{%
\section{Deep Dive: Forecast Distributions to Robust Decisions}
\begin{tcbitemize}[skin=sectionraster]
\tcbitem[title={Evaluate at the Decision Horizon},raster multicolumn=3]
Use rolling-origin evaluation with the forecast horizons, hierarchy, and data
latency of the planning process.
Report scale-appropriate errors and interval coverage; one aggregate point
metric can hide poor service for intermittent or high-value items.

\tcbitem[title={Reconcile the Hierarchy},raster multicolumn=3]
Product, region, channel, and total forecasts should add consistently.
Reconciliation trades local fit for coherent planning and must respect which
levels contain reliable signal and which are management aggregates.

\tcbitem[title={Optimise Across Scenarios},raster multicolumn=6]
Pass a forecast distribution or scenarios into inventory, routing, or capacity
decisions rather than treating a point forecast as certain demand.
Stress lead time, supplier failure, correlated disruption, and policy delay.
Compare the learned policy with a transparent robust or stochastic-optimisation
baseline.
\end{tcbitemize}
}

\expandafter\def\csname almanachdeepdive@multi-agent-systems\endcsname{%
\section{Deep Dive: LLM Agents as Distributed Systems}
\begin{tcbitemize}[skin=sectionraster]
\tcbitem[title={Distinguish Three Agent Traditions},raster multicolumn=3]
Classical BDI agents execute explicit beliefs, desires, and intentions.
MARL agents learn policies through interaction.
LLM agents generate plans and tool calls from context.
The three can be combined, but their state, guarantees, and failure modes are
not interchangeable.

\tcbitem[title={Protocols Need Semantics},raster multicolumn=3]
A message schema states syntax; an interaction protocol also defines allowed
states, idempotency, timeouts, authorization, conflict resolution, and recovery.
Use correlation identifiers and explicit ownership so retries do not create
duplicate consequential actions.

\tcbitem[title={Evaluate Trajectories and Emergence},raster multicolumn=6]
Measure task success, steps, tool errors, communication cost, latency,
permission violations, deadlocks, loops, and human interventions.
Test adversarial messages and collusion as well as cooperative examples.
Local pass rates do not prove stable collective behaviour.
\end{tcbitemize}
}

\expandafter\def\csname almanachdeepdive@commerce-marketing\endcsname{%
\section{Deep Dive: Feedback Loops and Multi-Sided Objectives}
\begin{tcbitemize}[skin=sectionraster]
\tcbitem[title={Logged Data Reflects the Old Policy},raster multicolumn=3]
Clicks and purchases are observed only for items exposed by an earlier ranking
and interface.
Position, availability, and selection bias confound naive supervised targets.
Randomised exploration or carefully justified propensity methods help estimate
counterfactual performance.

\tcbitem[title={Delayed and Interfering Outcomes},raster multicolumn=3]
Campaigns can change later conversion, returns, fatigue, and word of mouth.
One user's exposure can affect inventory, prices, sellers, and other users, so
standard independent-unit A/B assumptions may fail.

\tcbitem[title={Balance the Marketplace},raster multicolumn=6]
Evaluate relevance, conversion, margin, diversity, novelty, long-term retention,
producer exposure, support burden, and safety constraints.
Define which objectives are hard constraints and which can trade off.
Monitor feedback loops that make the training data increasingly reflect the
model's own past choices.
\end{tcbitemize}
}

\expandafter\def\csname almanachdeepdive@it-law\endcsname{%
\section{Deep Dive: From Legal Classification to Engineering Evidence}
\begin{tcbitemize}[skin=sectionraster]
\tcbitem[title={Begin with Role and Intended Purpose},raster multicolumn=3]
Map jurisdictions, actor roles, intended purpose, affected people, sector,
data flows, and lifecycle stage before selecting obligations.
Provider, deployer, importer, distributor, employer, controller, and processor
can carry different duties for the same technical component.

\tcbitem[title={Translate Duties into Controls},raster multicolumn=3]
Turn requirements into owners, controls, evidence, retention, review dates, and
incident paths.
Examples include provenance records, risk management, human oversight,
cybersecurity, performance monitoring, notices, rights handling, and supplier
contracts.

\tcbitem[title={Maintain a Legal Change Log},raster multicolumn=6]
Record the source, jurisdiction, version, effective date, interpretation owner,
affected system, and next review.
The legal chapter is orientation, not legal advice; current official text and
qualified counsel govern consequential decisions.
Reassess when purpose, data, model, geography, users, or autonomy changes.
\end{tcbitemize}
}

\expandafter\def\csname almanachdeepdive@visualization\endcsname{%
\section{Deep Dive: Uncertainty, Accessibility, and Monitoring}
\begin{tcbitemize}[skin=sectionraster]
\tcbitem[title={Show the Distribution Behind the Mean},raster multicolumn=3]
Use intervals, densities, small multiples, or raw observations to show spread,
sample size, and heterogeneity.
For probabilistic models, reliability diagrams and risk--coverage curves reveal
behaviour hidden by a single discrimination score.

\tcbitem[title={Design Beyond Colour},raster multicolumn=3]
Use contrast, direct labels, position, shape, line style, and ordering so the
message survives colour-vision differences and monochrome output.
Provide a meaningful caption or alternate description and preserve logical
reading order in the exported document.

\tcbitem[title={A Monitoring Chart Must Trigger Action},raster multicolumn=6]
Connect every panel to an owner, threshold, investigation, and time horizon.
Separate service health, data quality, prediction distribution, delayed
performance, subgroup behaviour, incidents, and business outcomes.
Annotate model or policy changes so the chart can support causal investigation
rather than merely display motion.
\end{tcbitemize}
}

\expandafter\def\csname almanachdeepdive@startup\endcsname{%
\section{Deep Dive: Portfolio Evidence and Option Value}
\begin{tcbitemize}[skin=sectionraster]
\tcbitem[title={Fund Learning Milestones},raster multicolumn=3]
Tie spending to the next uncertainty-reducing milestone: verified problem,
representative data, baseline feasibility, paid workflow adoption, responsible
operation, or repeatable economics.
Do not confuse accumulated code with reduced venture risk.

\tcbitem[title={Preserve Strategic Options},raster multicolumn=3]
Design pilots so data, evaluations, interfaces, and customer learning remain
useful if the model, supplier, or route to market changes.
Portability and a credible manual fallback create option value when technology
and regulation move faster than the plan.
\end{tcbitemize}
}

\expandafter\def\csname almanachdeepdive@scientific-writing\endcsname{%
\section{Deep Dive: Reproducibility Audit}
\begin{tcbitemize}[skin=sectionraster]
\tcbitem[title={Recompute from a Clean Environment},raster multicolumn=3]
Give the repository and manuscript to a reader who did not create the original
environment.
Record missing data, secrets, undocumented preprocessing, manual steps,
unstable downloads, and discrepancies between figures and reported numbers.

\tcbitem[title={Audit the Claim Graph},raster multicolumn=3]
For every abstract and conclusion claim, link the supporting result, method,
artifact, and limitation.
Remove orphan claims and propagate changed numbers or caveats through captions,
tables, discussion, model cards, and supplementary material.
\end{tcbitemize}
}

\newcommand{\almanachdeepdive}[1]{%
    \ifcsname almanachdeepdive@#1\endcsname
        \csname almanachdeepdive@#1\endcsname
    \fi
}






% ----------------------------------------------------------------------------------------------------------------------
% Header and Footer Settings
% ----------------------------------------------------------------------------------------------------------------------

\fancyhead[L]{
    \LARGE {\headerfont AI Almanach} \textcolor{\maincolor}{for engineers}
}
\fancyhead[C]{}
\fancyhead[R]{}
\fancyfoot[L]{}
\fancyfoot[C]{}
\fancyfoot[R]{\thepage}


% ----------------------------------------------------------------------------------------------------------------------
% Main Document
% ----------------------------------------------------------------------------------------------------------------------


\begin{document}

    \pagenumbering{gobble}
    % Front matter for AI Almanach.
% Keep this file independent from chapter numbering.

\begin{titlepage}
    \thispagestyle{empty}
    \begin{tikzpicture}[remember picture,overlay]
        \fill[\maincolor!8!white]
            (current page.south west) rectangle (current page.north east);
        \fill[\maincolor!75!black]
            (current page.north west) rectangle
            ([yshift=-3.2cm]current page.north east);
        \draw[\maincolor!45!white,line width=1.2pt]
            ([xshift=1.3cm,yshift=-4.3cm]current page.north west)
            -- ([xshift=-1.3cm,yshift=-4.3cm]current page.north east);
        \foreach \x/\y/\r in {
            2.0/4.0/0.18, 4.1/5.0/0.13, 6.2/3.7/0.16,
            8.0/5.4/0.12, 10.2/4.2/0.19, 12.1/5.1/0.14,
            14.0/3.8/0.17, 16.2/5.0/0.13, 18.1/4.0/0.18}
        {
            \fill[\maincolor!60!black]
                ([xshift=\x cm,yshift=\y cm]current page.south west)
                circle (\r cm);
        }
        \draw[\maincolor!35!black,line width=0.6pt]
            ([xshift=2.0cm,yshift=4.0cm]current page.south west)
            -- ([xshift=4.1cm,yshift=5.0cm]current page.south west)
            -- ([xshift=6.2cm,yshift=3.7cm]current page.south west)
            -- ([xshift=8.0cm,yshift=5.4cm]current page.south west)
            -- ([xshift=10.2cm,yshift=4.2cm]current page.south west)
            -- ([xshift=12.1cm,yshift=5.1cm]current page.south west)
            -- ([xshift=14.0cm,yshift=3.8cm]current page.south west)
            -- ([xshift=16.2cm,yshift=5.0cm]current page.south west)
            -- ([xshift=18.1cm,yshift=4.0cm]current page.south west);
    \end{tikzpicture}

    \vspace*{1.1cm}
    {\color{white}\sffamily\bfseries\fontsize{34}{38}\selectfont
        AI Almanach\par}
    \vspace{0.35cm}
    {\color{white}\sffamily\fontsize{16}{20}\selectfont
        Machine Learning, Data Science, and Intelligent Systems\par}
    \vspace{0.15cm}
    {\color{white}\sffamily\fontsize{13}{16}\selectfont
        A step-by-step engineering reference\par}

    \vfill
    \begin{tcolorbox}[
        enhanced,
        width=0.72\textwidth,
        colback=white!94!\maincolor,
        colframe=\maincolor!70!black,
        boxrule=0.6pt,
        arc=1mm]
        \sffamily
        From vectors to validated AI systems; from one gradient step to a
        tiny language model with a LoRA adapter.
        The destination is competence, not vocabulary bingo.
    \end{tcolorbox}
    \vfill

    {\sffamily\Large Bernhard Gerlach\par}
    \vspace{0.25cm}
    {\sffamily\large AI Almanach, version 2 working edition\par}
    {\sffamily\normalsize \today\par}
\end{titlepage}

\pagenumbering{roman}
\clearpage
\thispagestyle{empty}
\vspace*{\fill}
{\sffamily\small
\textbf{Purpose.}
This AI Almanach is an educational engineering reference, not a substitute for
domain-specific professional advice, safety assurance, or legal review.

\medskip
\textbf{Scientific status.}
Definitions, equations, algorithms, and empirical claims should be traceable
to cited primary literature, standards, official documentation, or clearly
identified secondary sources.
Claims about fast-moving models, software, law, and benchmarks carry a source
date and must be rechecked before consequential use.

\medskip
\textbf{Reproducibility.}
Practical results depend on data, random seeds, software versions, numerical
precision, hardware, and measurement protocol.
A benchmark without its conditions is a story, not a measurement.

\medskip
Typeset with LuaLaTeX using a card-based design built on \texttt{tcolorbox},
\texttt{TikZ}, and \texttt{biblatex}.
}
\vspace*{\fill}
\clearpage

\chapter*{Abstract}
\addcontentsline{toc}{chapter}{Abstract}

Artificial intelligence is not one technique.
It is a family of ways to represent problems, learn from data, reason under
uncertainty, optimise decisions, and build systems that interact with people
and the physical world.
The AI Almanach develops that family from first principles for engineering
students with mixed backgrounds.

The route begins with mathematical and programming refreshers, continues
through data analysis and classical machine learning, and then builds neural
networks, reinforcement learning, computer vision, natural language
processing, generative models, large language models, and multi-agent systems.
Applications are connected to the engineering constraints that determine
whether a model is useful: data quality, evaluation design, latency, memory,
energy, safety, security, privacy, maintainability, and human oversight.

The central didactic rule is simple: one concept at a time, but never in
isolation.
Every important idea is introduced through motivation, plain-language
intuition, notation and assumptions, a worked example, an engineering note,
a failure mode, a check for understanding, and a practical application.
Testing and evaluation form a continuous spine rather than a final chapter
that can be politely ignored when a deadline approaches.

\chapter*{Preface: learning AI without losing the plot}
\addcontentsline{toc}{chapter}{Preface}

AI can feel difficult for an innocent reason: several kinds of knowledge are
stacked on top of one another.
A learner may be meeting a new application, a new mathematical notation, a
new algorithm, and a new software tool in the same paragraph.
That is four lessons wearing one trench coat.

This book separates those layers and then reconnects them deliberately.
You will repeatedly move through the same engineering loop:

\begin{center}
\begin{tikzpicture}[
    node distance=8mm and 10mm,
    every node/.style={
        draw=\maincolor!65!black,
        fill=\maincolor!5!white,
        rounded corners=1mm,
        font=\sffamily\small,
        align=center,
        minimum height=8mm},
    >={Stealth[length=2mm]}]
    \node (question) {Question\\and constraints};
    \node[right=of question] (data) {Data and\\assumptions};
    \node[right=of data] (model) {Model and\\objective};
    \node[below=of model] (eval) {Evaluation\\and uncertainty};
    \node[left=of eval] (system) {System and\\operations};
    \node[left=of system] (decision) {Decision, feedback,\\and revision};
    \draw[->] (question) -- (data);
    \draw[->] (data) -- (model);
    \draw[->] (model) -- (eval);
    \draw[->] (eval) -- (system);
    \draw[->] (system) -- (decision);
    \draw[->] (decision) -- (question);
\end{tikzpicture}
\captionof{figure}{The engineering learning loop used throughout the AI Almanach.}
\label{fig:engineering-learning-loop}
\end{center}

The loop matters because model training is only one step.
An elegant model trained on a leaked target, evaluated with the wrong metric,
or deployed without monitoring is still an unsuccessful engineering system.

\chapter*{How to use the AI Almanach}
\addcontentsline{toc}{chapter}{How to use the AI Almanach}

\section*{Choose a route}

The chapters are dependency ordered, but the book is also a reference.
Choose the route that matches your immediate goal and return to prerequisites
when a symbol or assumption feels unfamiliar.

\begin{table}[htbp]
\centering
\caption{Suggested learning routes through the AI Almanach.}
\label{tab:learning-routes}
\begin{tblr}{
    width=\textwidth,
    colspec={Q[l,0.18\textwidth] X[l] X[l]},
    row{1}={font=\bfseries, bg=\maincolor!15!white},
    hlines,
    vlines}
Route & Begin with & Then emphasise \\
First complete pass & Orientation, mathematics, Python, and data &
Classical ML, deep learning, evaluation, then the domain chapters \\
ML engineer & Mathematics refresh and data foundations &
Classical ML, deep learning, evaluation/MLOps, and practical labs \\
LLM engineer & Linear algebra, probability, neural networks &
NLP, Transformers, practical LLM engineering, evaluation, safety \\
Data scientist & Statistics, Python, and data analytics &
Classical ML, causal reasoning, visualisation, and deployment \\
Technical leader & Introduction and evaluation &
Safety, law, operations, applications, and evidence limitations \\
\end{tblr}
\end{table}

\section*{The standard learning unit}

New and substantially revised material follows the same sequence.
The sequence reduces cognitive switching and makes omissions visible during
review.

\begin{enumerate}
    \item \textbf{Why it matters:} the question the concept helps answer.
    \item \textbf{Prerequisites:} the minimum earlier ideas and notation.
    \item \textbf{Plain language:} a useful mental model, including its limits.
    \item \textbf{Formal model:} definitions, shapes, units, and assumptions.
    \item \textbf{Worked example:} small enough to calculate by hand.
    \item \textbf{Engineer's note:} implementation and systems consequences.
    \item \textbf{Failure mode:} how the idea is commonly misused.
    \item \textbf{Check:} a retrieval or transfer question with a testable
          answer.
    \item \textbf{Practice:} a reproducible experiment with a baseline.
    \item \textbf{Evidence:} citations and the scope of what they support.
\end{enumerate}

\begin{engineernote}
You do not understand an algorithm merely because its name feels familiar.
You understand it when you can state its inputs and outputs, draw the data
flow, calculate a toy case, identify its assumptions, predict a failure mode,
and design a test that could prove your implementation wrong.
\end{engineernote}

\section*{Why the learning pattern is repeated}

The repeated chapter compass, worked examples, failure cards, and retrieval
prompts are deliberate rather than decorative.
Constructive alignment connects observable learning outcomes to practice and
assessment\textsuperscript{\cite{biggsEnhancingTeachingConstructive1996}}.
Worked examples reduce avoidable search for novices, while support should fade
as expertise grows\textsuperscript{\cite{swellerCognitiveLoadTheory2019}}.
Practice testing and distributed practice have broad evidence for improving
retention\textsuperscript{\cite{dunloskyImprovingStudentsLearning2013}}, and
retrieval can outperform additional study after a delay
\textsuperscript{\cite{roedigerTestEnhancedLearning2006}}.

Words and explanatory graphics are placed together when the visual carries a
real part of the explanation.
This follows multimedia, signalling, coherence, and spatial-contiguity
principles; irrelevant illustration is omitted because visual density can add
extraneous processing instead of understanding
\textsuperscript{\cite{mayerUsingMultimediaElearning2017}}.

\begin{table}[htbp]
\centering
\caption{Evidence-informed learning patterns used throughout the book.}
\label{tab:evidence-informed-learning-patterns}
\begin{tblr}{
    width=\textwidth,
    colspec={Q[l,0.2\textwidth]X[l]X[l]},
    row{1}={font=\bfseries,bg=\maincolor!15!white},
    hlines,
    vlines}
Pattern & Purpose & How to use it \\
Chapter compass & Activate prerequisites and expose relationships & Predict
the four-stage chain before reading; redraw it afterwards \\
Worked example & Show the reasoning path, not only the answer & Explain why
each step is valid, then complete a partially worked variant \\
Failure card & Build diagnostic knowledge & Predict the symptom, locate the
violated assumption, and propose a discriminating test \\
Retrieval prompt & Strengthen later recall and expose gaps & Answer without
looking, check with feedback, and revisit after a delay \\
Transfer task & Prevent knowledge from remaining example-bound & Change the
domain, constraints, or data-generating process and adapt the method \\
Adjacent visual & Reduce split attention for spatial or process ideas & Trace
the diagram while explaining the corresponding text in your own words \\
\end{tblr}
\end{table}

These effects depend on prior knowledge, task complexity, feedback, and study
delay.
The book therefore offers a route rather than claiming that one layout works
identically for every reader.
Experienced readers may skip elementary scaffolds, but should still attempt
the failure and transfer prompts.

\section*{Evidence and verification policy}

The bibliography is not decoration.
Each citation should be close to the claim it supports, and important claims
should be triangulated where practical.

\begin{description}
    \item[Level A --- primary or normative evidence:]
        peer-reviewed papers, original technical reports, standards, laws,
        and official datasets or model cards.
    \item[Level B --- authoritative synthesis:]
        established textbooks, systematic reviews, and recognised handbooks.
    \item[Level C --- operational evidence:]
        official software documentation and reproducible implementation notes.
    \item[Level D --- illustration:]
        tutorials, analogies, and examples that clarify an idea but do not
        establish a general scientific claim.
\end{description}

For equations, verification means more than finding the same symbols in two
books.
Dimensions, units, boundary cases, sign conventions, and a numerical toy
example must agree.
For empirical claims, the dataset, split, metric, uncertainty, hardware, and
software configuration belong to the claim.

\begin{cautionbox}[The AI Almanach is a map, not the territory]
AI changes quickly and application domains contain specialised knowledge that
no single volume can exhaust.
Use the AI Almanach to build durable mental models and verification habits, then
read the cited primary sources for decisions that matter.
\end{cautionbox}

\section*{Notation quick reference}

\begin{table}[htbp]
\centering
\caption{Core notation used across chapters.}
\label{tab:core-notation}
\begin{tblr}{
    width=\textwidth,
    colspec={Q[c,0.13\textwidth] X[l] X[l]},
    row{1}={font=\bfseries, bg=\maincolor!15!white},
    hlines,
    vlines}
Symbol & Meaning & Typical shape or range \\
$x$ or $\mathbf{x}$ & scalar or feature vector &
$\mathbf{x}\in\mathbb{R}^{d}$ \\
$X$ & design matrix or token representations & $X\in\mathbb{R}^{n\times d}$ \\
$y$ & observed target & scalar, class, sequence, or structured object \\
$\hat y$ & model prediction & same semantic space as the target \\
$\theta$ & trainable model parameters & vector or collection of tensors \\
$\mathcal{L}$ & loss or empirical objective & usually $\mathbb{R}$ \\
$p(y\mid x)$ & conditional probability & $[0,1]$, sums or integrates to one \\
$\nabla_{\theta}\mathcal{L}$ & gradient of loss with respect to parameters &
same shape as $\theta$ \\
$\mathbb{E}[X]$ & expectation of a random variable & units of $X$ \\
\end{tblr}
\end{table}

\begin{checkpoint}
Choose a model you already know.
Can you name its input, target, parameters, objective, evaluation unit, one
assumption, and one failure mode?
Any blank is a useful reading goal rather than a reason for panic.
\end{checkpoint}


    \setcounter{tocdepth}{1}
    \tableofcontents
    \clearpage
    \listoffigures
    \clearpage
    \listoftables
    \clearpage

    \pagenumbering{arabic}

    % ==== Chapters (each in separate file) ========================================================
    \part{Orientation and Foundations}
    \almanachinputchapter{orientation}
        {How do intelligent systems turn goals and data into defensible
         decisions?}
        {Problem and context}{Representation}{Learning or reasoning}
        {Evaluation and action}{chapters/ch01-introduction}
    \almanachinputchapter{mathematics}
        {Which mathematical objects let us describe learning, uncertainty,
         and optimisation precisely?}
        {Objects and notation}{Transformations}{Uncertainty}
        {Optimisation}{chapters/ch02-mathematical-foundations}
    \almanachinputchapter{programming}
        {How do we turn an analytical idea into software that can be tested,
         reproduced, and maintained?}
        {Types and interfaces}{Arrays and data frames}{Pipelines}
        {Tests and versions}{chapters/ch03-programming-foundations}
    \almanachinputchapter{data-analytics}
        {How does raw data become trustworthy evidence for an engineering
         decision?}
        {Acquire and define}{Clean and integrate}{Explore and model}
        {Govern and communicate}{chapters/ch04-data-analytics}

    \part{Learning from Data}
    \almanachinputchapter{machine-learning}
        {How can a model learn patterns that generalise beyond the examples
         used to fit it?}
        {Task and baseline}{Representation and model}{Fit and regularise}
        {Validate and diagnose}{chapters/ch05-machine-learning}
    \almanachinputchapter{deep-learning}
        {How do layered differentiable models learn useful representations at
         scale?}
        {Represent}{Differentiate}{Optimise}{Scale and verify}
        {chapters/ch06-deep-learning}
    \almanachinputchapter{evaluation-mlops}
        {What evidence shows that a learned system is useful now and remains
         useful after deployment?}
        {Evaluation contract}{Tests and uncertainty}{Release controls}
        {Monitoring and revision}{chapters/ch06b-evaluation-mlops}
    \almanachinputchapter{reinforcement-learning}
        {How should an agent learn sequential decisions when actions change
         the data it will observe next?}
        {State and objective}{Act and observe}{Update value or policy}
        {Evaluate behaviour}{chapters/ch07-reinforcement-learning}

    \part{Trustworthy Perception, Language, and Generation}
    \almanachinputchapter{trustworthy-ai}
        {What evidence justifies trusting an AI-enabled system within a stated
         operating boundary?}
        {Scope and harms}{Threats and hazards}{Controls and tests}
        {Assurance and operation}{chapters/ch08-functional-security}
    \almanachinputchapter{computer-vision}
        {How do pixels and sensor geometry become useful, testable statements
         about a scene?}
        {Image formation}{Features and geometry}{Prediction}
        {Robust perception}{chapters/ch09-computer-vision}
    \almanachinputchapter{nlp-generative-ai}
        {How do language systems represent context and generate outputs that
         can be evaluated rather than merely admired?}
        {Tokenise}{Represent context}{Generate or retrieve}
        {Evaluate meaning and risk}{chapters/ch12-nlp-voice-assistants}
    \almanachinputchapter{practical-llm-engineering}
        {How can a small language-model experiment reveal the mechanics and
         trade-offs of real LLM systems?}
        {Trace one token}{Adapt weights}{Benchmark quality and cost}
        {Package and deploy}{chapters/ch12b-practical-llm-engineering}

    \part{AI in Engineering and Society}
    \almanachinputchapter{fintech}
        {How can AI support financial decisions while controlling asymmetric
         error costs, abuse, and regulatory risk?}
        {Decision and data}{Risk model}{Constraint and control}
        {Audit and monitoring}{chapters/ch10-fintech}
    \almanachinputchapter{healthcare}
        {How can AI improve a clinical workflow without confusing benchmark
         performance with patient benefit?}
        {Clinical need}{Representative evidence}{Workflow integration}
        {Outcome monitoring}{chapters/ch11-healthcare}
    \almanachinputchapter{industrial-ai}
        {How do predictions become safe, timely actions in a physical
         production system?}
        {Sense}{Predict}{Plan}{Control and recover}
        {chapters/ch13-industrial-ai}
    \almanachinputchapter{supply-chain}
        {How can uncertain demand and constrained resources be coordinated
         across a changing network?}
        {Observe the network}{Forecast scenarios}{Optimise decisions}
        {Monitor and recover}{chapters/ch14-supply-chain}
    \almanachinputchapter{multi-agent-systems}
        {How can autonomous components coordinate without assuming that local
         intelligence creates global order?}
        {Perceive and reason}{Communicate}{Coordinate incentives}
        {Verify emergence}{chapters/ch15-multi-agent-systems}
    \almanachinputchapter{commerce-marketing}
        {How can personalisation and forecasting create value without turning
         feedback loops into manipulation or bias?}
        {Observe behaviour}{Predict and personalise}{Experiment}
        {Govern outcomes}{chapters/ch16-ecommerce-marketing}
    \almanachinputchapter{it-law}
        {Which legal duties attach to an AI system, its data, its outputs, and
         the organisations operating it?}
        {Classify the use}{Identify duties}{Document evidence}
        {Reassess change}{chapters/ch17-it-law}
    \almanachinputchapter{visualization}
        {How can a visual representation make evidence easier to reason about
         without overstating what the data supports?}
        {Question}{Encode}{Compare}{Explain and verify}
        {chapters/ch18-visualization}
    \almanachinputchapter{startup}
        {How can an AI venture learn whether a painful problem, a feasible
         system, and sustainable economics exist at the same time?}
        {Problem evidence}{Small experiment}{Unit economics}
        {Scale responsibly}{chapters/ch19-startup}
    \almanachinputchapter{scientific-writing}
        {How can a technical claim be made traceable, reproducible, and useful
         to a sceptical reader?}
        {Claim}{Method}{Evidence}{Revision and publication}
        {chapters/ch20-scientific-writing}

    % ==== References ==============================================================================
    \part{Back Matter}
    \appendix
\chapter{Curriculum and Course Coverage Map}
\label{ch:course-coverage-map}

This appendix makes the curriculum auditable.
It maps the supplied external machine-learning course outcomes to the AI Almanach
and adds the broader competencies required to build dependable AI systems.

\begin{table}[htbp]
\centering
\caption{Coverage of the supplied external machine-learning course.}
\label{tab:external-course-coverage}
\begin{tblr}{
    width=\textwidth,
    colspec={X[l] Q[l,0.27\textwidth] X[l]},
    row{1}={font=\bfseries, bg=\maincolor!15!white},
    hlines,
    vlines}
Course requirement & Primary location & Evidence of mastery \\
Statistics, derivatives, and matrices for ML &
Mathematical Foundations &
Derive a gradient, check tensor dimensions, and quantify uncertainty \\
Classification and regression &
Machine Learning &
Implement baselines and compare metrics on held-out data \\
Model optimisation &
Deep Learning; Evaluation and MLOps &
Tune only on validation data and report the search protocol \\
Decision trees and forests &
Machine Learning &
Explain a split, control complexity, and measure generalisation \\
Neural networks and support vector machines &
Machine Learning; Deep Learning &
Calculate a toy forward pass and diagnose optimisation behaviour \\
Clustering and recommender systems &
Machine Learning; E-Commerce and Marketing &
Evaluate without pretending that unlabeled structure has one true answer \\
Develop, implement, and scientifically evaluate an ML approach &
Evaluation and MLOps; Practical LLM Engineering &
Reproduce a baseline-to-candidate experiment with uncertainty and error
analysis \\
\end{tblr}
\end{table}

\section{Competencies added by the AI Almanach}

\begin{itemize}
    \item formulate a decision problem before selecting a model;
    \item create data sheets, model cards, and reproducible experiment records;
    \item detect leakage, distribution shift, shortcut learning, and metric
          mismatch;
    \item measure latency, throughput, memory, energy, and operational
          reliability;
    \item evaluate generative outputs with task-specific, human, and safety
          tests;
    \item reason about privacy, security, functional safety, fairness, and law;
    \item communicate evidence and uncertainty to technical and non-technical
          readers.
\end{itemize}

\section{Alignment with current curriculum frameworks}

The curriculum also uses the ACM/IEEE/AAAI computing competencies as a
technical breadth check and UNESCO's AI competency framework as a human-centred
literacy check\textsuperscript{\cite{acmCS2023Curricula2024,unescoAICompetencyFramework2024}}.
These frameworks are complementary rather than interchangeable: one helps
audit computing depth, while the other foregrounds human agency, ethics, AI
techniques, and system design.

\begin{table}[htbp]
\centering
\caption{Framework alignment used to audit the AI Almanach curriculum.}
\label{tab:curriculum-framework-alignment}
\begin{tblr}{
    width=\textwidth,
    colspec={Q[l,0.22\textwidth] X[l] X[l]},
    row{1}={font=\bfseries, bg=\maincolor!15!white},
    hlines,
    vlines}
Audit lens & Almanach coverage & Observable evidence \\
Technical foundations &
Mathematics, programming, data, ML, deep learning, RL, vision, and NLP &
Trace a mechanism, implement a baseline, and test generalisation \\
Systems and operations &
Evaluation, MLOps, LLM engineering, security, and industrial systems &
Specify release gates, operational budgets, monitoring, and rollback \\
Human agency and ethics &
Trustworthy AI, healthcare, finance, commerce, law, and visualization &
Identify affected actors, contestability, harms, uncertainty, and duties \\
Design and creation &
Domain chapters, multi-agent systems, venture building, and scientific
writing &
Produce a bounded solution, evidence record, and reproducible argument \\
\end{tblr}
\end{table}

\section{Definition of done for a chapter}

A chapter is considered reviewed only when its important concepts have
prerequisites, definitions, dimensional checks, at least one worked example,
failure modes, evaluation guidance, and traceable sources.
The chapter must compile without hidden local assets and must not contain
unresolved placeholders.
Time-sensitive claims must include a verification date or be rewritten as
stable principles.

    \chapter{References}
    \printbibliography[heading=none]

\end{document}
